%----------------------------------------------------------------------------------------
% Tailored CV — Wissenschaftlicher Mitarbeiter, Genetische Epidemiologie / Statistische Genetik
% Institut für Medizinische Statistik, UKSH Kiel — Kennziffer 28393
%----------------------------------------------------------------------------------------

\documentclass[10pt,a4paper,sans]{moderncv}

\usepackage[utf8]{inputenc}
\usepackage[T1]{fontenc}

\moderncvstyle{classic}
\moderncvcolor{blue}

\usepackage[scale=0.78]{geometry}
\usepackage{ragged2e}

%----------------------------------------------------------------------------------------

\firstname{Kundai Farai}
\familyname{Sachikonye}

\address{Auf der Lichtung 19}{82178 Puchheim, Germany}
\mobile{(+49) 1626386344}
\email{kundai.sachikonye@bitspark.com}
\homepage{github.com/fullscreen-triangle}{github.com/fullscreen-triangle}

%----------------------------------------------------------------------------------------

\begin{document}

\makecvtitle

%----------------------------------------------------------------------------------------
%	EDUCATION
%----------------------------------------------------------------------------------------

\section{Education}

\cventry{2019--2021}{PhD candidate, Computational Lipidomics}{Technische Universität München / Bayerisches Forschungsinstitut}{}{}{
  Multi-omics statistical modelling, mass spectrometry pipelines, federated
  learning across distributed clinical sites. Departed to pursue independent research.
}
\cventry{2018--2019}{MSc Computational Biology}{Universität Tübingen}{}{}{
  Statistical genetics, bioinformatics, sequence analysis, variant calling,
  RNA-seq quantification, population genetics.
}
\cventry{2014--2017}{MSc Pharmaceutical Biotechnology}{HAW Hamburg}{}{}{}
\cventry{2011--2014}{BSc Biochemistry and Cell Biology}{Jacobs University Bremen}{}{}{}

%----------------------------------------------------------------------------------------
%	RELEVANT PUBLICATIONS
%----------------------------------------------------------------------------------------

\section{Relevant Manuscripts}

\cvitem{2025}{
  \textbf{On the Thermodynamic Consequences of Categorical Completion on Biological Systems.}
  MHC binding affinity correlates with categorical richness ($r=0.73$, $p<10^{-12}$,
  $n=2{,}847$ IEDB peptides). AUC 0.81 vs NetMHCpan 0.68.
  \textit{Manuscript.}
}

\cvitem{2025}{
  \textbf{On Disease Epistemology in Blind Receiver.}
  Loop-holonomy self-consistency diagnostic; AUC$=1.00$; 25/25 validation experiments;
  sparse $\ell_1$ LP for personalised therapeutic target selection.
  \textit{Manuscript. AIMe Registry / TUM Weihenstephan.}
}

\cvitem{2025}{
  \textbf{Syndrome: Categorical Resolution of Disease Trajectories.}
  Eight-dimensional disease vector $\mathbf{D}\in[0,1]^8$; $\mathcal{O}(N\log N)$
  trajectory computation; minimum-intervention LP for personalised risk management.
  \textit{Manuscript.}
}

%----------------------------------------------------------------------------------------
%	RESEARCH SOFTWARE
%----------------------------------------------------------------------------------------

\section{Research Software}

\cvitem{gospel}{
  \textbf{Genome-Wide Variant Analysis and Multi-Omics Integration.}
  $10^6$ variants/second VCF processing; RNA-seq quantification integration;
  Bayesian network inference and Gaussian mixture models for multi-omics
  statistical modelling. Validated: 1000~Genomes Phase~3, gnomAD~v3.1.2,
  UK~Biobank. Pathway analysis (KEGG, Reactome); protein structure (RDKit).
  Applicable to polygenic score construction and familial genomic risk stratification.
  \href{https://github.com/fullscreen-triangle/gospel}{github.com/fullscreen-triangle/gospel}
}

\cvitem{syndrome}{
  \textbf{Individual-Level Disease Risk Prediction and Trajectory.}
  Disease state vector $\mathbf{D}\in[0,1]^8$ across eight pathophysiological
  dimensions (immune activation, barrier integrity, genetic load, epigenetic
  state, etc.); per-dimension coherence indices $\eta_i\in[0,1]$ from
  biomarker time-series. $\mathcal{O}(N\log N)$ trajectory to healthy attractor;
  sparse $\ell_1$ LP for minimum-intervention therapeutic target selection.
  Loop-holonomy self-consistency: AUC$=1.00$, 25/25 experiments.
  Extends polygenic risk scores by integrating genetic risk load with
  non-genetic pathophysiological components into a structured risk coordinate.
  \href{https://github.com/fullscreen-triangle/syndrome}{github.com/fullscreen-triangle/syndrome}
}

\cvitem{lavoisier}{
  \textbf{Large-Scale Multi-Omics Data Pipeline.}
  Dual-pipeline spectral analysis: numerical + CNN (PyTorch); 98.9\,\% NIST
  accuracy; Ray distributed; $>$100\,GB dataset support; memory-mapped I/O.
  \href{https://github.com/fullscreen-triangle/lavoisier}{github.com/fullscreen-triangle/lavoisier}
}

%----------------------------------------------------------------------------------------
%	EXPERIENCE
%----------------------------------------------------------------------------------------

\section{Experience}

\cventry{2021--present}{Independent AI \& Computational Biology Developer}{Bitspark GmbH}{}{}{
  Genomic analysis frameworks, multi-omics integration, disease risk prediction
  models, statistical methods for high-dimensional biological data. All systems
  publicly documented and validated.
}

\cventry{2019--2021}{Computational Lipidomics}{TUM / Bayerisches Forschungsinstitut}{Munich}{}{
  Mass spectrometry-based lipidomics: peak detection, ion annotation (LIPID MAPS,
  LipidBlast), multi-omics statistical modelling, federated learning for
  distributed clinical analysis. Python, Java.
}

\cventry{2017--2018}{Glycomics and Proteomics}{Universitätsklinikum Hamburg-Eppendorf}{}{}{
  Automated LC-MS/MS and MALDI-TOF pipeline development; sample preparation;
  reproducibility validation; SOP development.
}

\cventry{2013}{Cancer Biomarker Discovery}{University Medical Center Groningen}{}{}{
  UPLC-MS/MS shotgun proteomics for cancer biomarker discovery; statistical
  validation against reference standards.
}

%----------------------------------------------------------------------------------------
%	TECHNICAL SKILLS
%----------------------------------------------------------------------------------------

\section{Technical Skills}

\cvitem{Statistical Genetics}{
  Genome-wide variant analysis (VCF, PLINK format); population structure
  (PCA, admixture); polygenic score methodology (C+T, Bayesian shrinkage,
  continuous shrinkage priors); Bayesian networks for genomic co-occurrence;
  linkage disequilibrium modelling; familial/pedigree data structures;
  mixed-effects models for correlated observations
}

\cvitem{Statistical Methods}{
  Bayesian networks; variational Bayes; Gaussian mixture models;
  sparse $\ell_1$ regularisation (LP); graph Laplacian methods;
  Banach contraction / fixed-point theory; survival analysis; GLMs
}

\cvitem{Languages}{
  R (primary for statistical modelling), Python (data processing, ML pipelines),
  Rust, Bash, SQL
}

\cvitem{Genomic Data}{
  VCF, FASTQ, BAM; PLINK bed/bim/fam; gnomAD, 1000~Genomes, UK~Biobank;
  KEGG, Reactome, GSEA pathway analysis; LIPID MAPS, HMDB
}

\cvitem{Infrastructure}{
  Ray, Dask, Docker, FastAPI; memory-mapped I/O for large datasets;
  HDF5, mzML; Git
}

%----------------------------------------------------------------------------------------
%	LANGUAGES
%----------------------------------------------------------------------------------------

\section{Languages}

\cvitemwithcomment{English}{Native / Fluent}{}
\cvitemwithcomment{German}{Conversational}{Living in Germany since 2011}

%----------------------------------------------------------------------------------------

\renewcommand{\listitemsymbol}{-~}

\end{document}
